\chapter{Existing Systems/State of the Art}
\section{Overview of the Field}
Question answering over a specialist corpus involves two related but separate tasks: locating relevant information and producing an explanation that follows from it. Traditional retrieval exposes ranked documents or passages directly. Generative systems add a natural-language synthesis step, which can help users read the evidence but also introduces statements that require verification. For legal material, document numbers, provision boundaries, and qualification clauses make this distinction particularly consequential.

The present project belongs to retrieval-augmented question answering. Its comparison is internal: the same hybrid retrieval baseline is augmented with optional reranking, graph traversal, and feedback. This report reviews the methods that motivate those additions without transferring published performance claims to the Vietnamese development corpus.

\section{Historical Context}
Lexical retrieval, represented here by BM25, ranks documents using term evidence and length normalization \citep{robertson2009bm25}. It is useful when a query contains exact legal names or identifiers. Dense retrieval represents text in a learned vector space and can match related wording even when terms differ. Combining the two is a reasonable design hypothesis because exact and semantic matching can fail in different ways.

RAG introduced a prominent framework for conditioning language generation on retrieved non-parametric information \citep{lewis2020rag}. The original research includes model training and experiments on knowledge-intensive NLP tasks. This project borrows the retrieval-conditioned design principle while using frozen pretrained inference models; it does not reproduce the original training procedure or evaluation datasets.

\section{Key Studies and Theories}
\subsection{Retrieval Fusion and Ranking}
RRF combines ranked lists by summing reciprocal rank contributions \citep{cormack2009rrf}. Its appeal here is that BM25 scores, cosine similarities, and graph-derived ranks need not be placed on a common numerical scale. The method rewards candidates supported by several lists, while retaining contributions from a single retriever. This is a motivation for its use, not a guarantee of a gain for every corpus.

Ranking should also be evaluated in more than one way. Recall measures whether labeled evidence is found, reciprocal rank emphasizes the first relevant hit, and average precision accounts for all labeled hits. nDCG discounts relevance lower in the list and normalizes against an ideal ranking \citep{jarvelin2002gain}. The project uses binary chunk relevance, so it does not evaluate graded legal authority or graded degrees of usefulness.

\subsection{Multilingual Representations and Legal Structure}
M3-Embedding supports multilingual text representations and several retrieval modes \citep{chen2024m3}. The project uses only the dense representation of BGE-M3. Its separate BM25 component is not the learned sparse mode of that model. Therefore, calling the implementation ``BGE-M3 hybrid retrieval'' would obscure the actual combination being tested.

The cross-encoder model card describes a multilingual mMARCO-trained query--passage ranker \citep{mmarco}. A cross-encoder evaluates a question and candidate jointly rather than relying only on independently stored vectors. The implementation runs a pinned ONNX int8 export. Published multilingual capabilities do not establish Vietnamese legal-domain calibration, especially for long provisions truncated to the configured 512-token reranker input.

Legal hierarchy motivates a complementary representation choice. A leaf may be appropriate for retrieval, while a complete clause or article is needed for answering. This project explicitly separates embedded leaves from unembedded parents. The current experiment fixes that representation for every method; it provides no causal estimate that hierarchical chunking is superior to fixed windows.

\subsection{Corpus Provenance}
The collection is obtained from the \code{tmquan/phapdien-moj-gov-vn} dataset mirror \citep{phapdien}. Ingestion records the resolved source commit and preserves document and provision metadata. The benchmark uses a selected three-topic corpus, not the entire dataset. A provenance record makes the experiment repeatable at a given snapshot, but does not certify that the mirror is complete, legally reviewed, or up to date.

\section{Technological Advancements}
\subsection{Graph-Based Retrieval}
Graph approaches can expose relationships that a passage-only ranking may not directly surface. However, ``Graph RAG'' covers different architectures. \citet{edge2024graph} describe an LLM-built entity graph with community summaries for global, query-focused summarization. The implementation in this project instead uses explicit document references and source-backed typed links, plus a hierarchy of legal chunks. It has no community detection, community summaries, or global map--reduce synthesis.

This distinction determines the experiment's interpretation. Enabling the graph branch tests whether bounded document-neighbour expansion adds useful candidates. It does not reproduce Microsoft GraphRAG or test its published global sensemaking setting. The pilot also needs questions that actually depend on available graph links before it can substantiate a graph-specific benefit.

\subsection{Context and Answer Evaluation}
RAGAs motivates evaluating different stages of a RAG pipeline rather than relying on a single answer score \citep{es2024ragas}. This project implements its own Ragas-style context and answer measures with stored judgment details. It does not call the RAGAs package, and numerical values should not be assumed interchangeable with a particular library version or prompt set.

ARES similarly treats evaluation as a measurement problem, using learned judges and human annotation to improve estimation \citep{saadfalcon2024ares}. Its relevance here is methodological: a model judge needs calibration and validation. The project does not implement ARES or its training procedure. Its own recorded judge disagreements make independent human review a necessary next step.

\section{Comparison of Existing Systems}
Table~\ref{tab:related} compares architectural capabilities. It is not an empirical leaderboard, and does not assign scores to untested public or commercial systems.

\begin{table}[H]\centering\small
\caption{Conceptual comparison of retrieval and answering approaches.}\label{tab:related}
\begin{tabularx}{\textwidth}{P{2.5cm}YYY}\toprule
\textbf{Approach} & \textbf{Evidence access} & \textbf{Typical strength} & \textbf{Limitation relevant here}\\\midrule
Lexical retrieval & Term-matched documents/passages & Exact identifiers and terminology & Wording mismatch; no answer synthesis\\
Dense retrieval & Nearest vector representations & Semantic matching & Can blur exact scope and identifiers\\
Hybrid RAG & Fused candidates and generated answer & Combines lexical and semantic evidence & Generation may exceed evidence\\
Community GraphRAG & Entity graph and community summaries & Corpus-level thematic synthesis & Different task and graph construction\\
This prototype & Hybrid ranks; optional document graph; packed provisions & Traceable local evidence and observable gates & Pilot-limited coverage and uncalibrated judge\\\bottomrule
\end{tabularx}
\end{table}

\section{Gaps in Literature/Technology}
General retrieval and RAG research does not determine how well a particular legal snapshot, tokenizer, model, and prompt combination will work. The project therefore addresses an applied evidence gap: a reproducible comparison of these components on a Vietnamese legal pilot with explicit source labels. A second gap is operational observability. Ranking improvements can disappear when context packing drops a required condition or when a sufficiency gate rejects an otherwise answerable question.

The current evidence exposes a further measurement gap. A semantic judge can penalize a faithful scope statement because it expects details absent from the question. Such errors are not resolved by collecting more model-generated scores alone. The benchmark must retain inspectable judgments and eventually compare them with independently adjudicated legal evidence.

\section{Justification for the Project}
The chosen architecture makes these questions testable within a course project. Lexical and vector retrieval provide a shared baseline. Graph retrieval and reranking can be switched independently. Context packing and online evaluation remain common, and all outputs are captured before offline scoring. This design supports incremental engineering and controlled ablation without claiming a new foundation model or a complete legal reasoning system.
